\documentclass[../../main.tex]{subfiles}

\begin{document}

\section{集合间的映射}\label{section:ch00-映射}

给定两个集合 $A$ 和 $B$, 从 $A$ 到 $B$ 的映射或函数意味着对 $A$ 的每个成员指派 $B$ 的一个成员给它. 在 $B$ 是实数集的情形下, 我们总是用“函数”这个词. 一般我们记这样的映射为 $f:A\to B$, 而对 $A$ 的每个成员 $x$, 我们记 $f(x)$ 为 $B$ 中指派给 $x$ 的成员. 对于 $A$ 的子集 $A'$, 我们定义
\begin{align*}
f(A')=\{b\mid b=f(a),\,a\text{ 为 }A'\text{ 的某个成员}\};
\end{align*}
$f(A')$ 称为 $A'$ 在 $f$ 下的象. 我们称集合 $A$ 为函数 $f$ 的定义域, 而称 $f(A)$ 为 $f$ 的象或值域. 若 $f(A)=B$, 函数 $f$ 称为是映上的. 若对 $f(A)$ 的每个成员 $b$ 恰有 $A$ 的一个成员 $a$ 使得 $b=f(a)$, 函数 $f$ 称为是一对一的. 既是一对一又是映上的映射 $f:A\to B$ 称为是可逆的, 我们说该映射建立了集合 $A$ 与 $B$ 之间的一一对应. 给定一个可逆映射 $f:A\to B$, 对 $B$ 中的每个点 $b$, 恰好存在 $A$ 中的一个成员 $a$ 使得 $f(a)=b$, 它被记为 $f^{-1}(b)$. 这个指派定义了映射 $f^{-1}:B\to A$, 称之为 $f$ 的逆. 两个集合 $A$ 和 $B$ 称为是对等的, 若存在从 $A$ 映到 $B$ 的可逆映射. 从集合论的观点看, 对等的两个集合是不可区分的.

给定两个映射 $f:A\to B$ 和 $g:C\to D$ 使得 $f(A)\subseteq C$, 则复合 $g\circ f:A\to D$ 定义为对每个 $x\in A$, $[g\circ f](x)=g(f(x))$. 不难看出可逆映射的复合是可逆的. 对于集合 $D$, 定义恒等映射 $\mathrm{id}_{D}:D\to D$ 为对所有 $x\in D$, $\mathrm{id}_{D}(x)=x$. 映射 $f:A\to B$ 是可逆的, 当且仅当存在映射 $g:B\to A$ 使得
\begin{align*}
g\circ f=\mathrm{id}_{A}\quad\text{且}\quad f\circ g=\mathrm{id}_{B}.
\end{align*}

即便映射 $f:A\to B$ 不是可逆的, 对于集合 $E$, 我们定义 $f^{-1}(E)$ 为集合 $\{a\in A\mid f(a)\in E\}$, 称之为 $E$ 在 $f$ 下的原象. 我们有下面有用的性质: 对于任何两个集合 $E_{1}$ 和 $E_{2}$,
\begin{align*}
f^{-1}(E_{1}\cup E_{2})=f^{-1}(E_{1})\cup f^{-1}(E_{2}),\qquad
f^{-1}(E_{1}\cap E_{2})=f^{-1}(E_{1})\cap f^{-1}(E_{2})
\end{align*}
与
\begin{align*}
f^{-1}(E_{1}\setminus E_{2})=f^{-1}(E_{1})\setminus f^{-1}(E_{2}).
\end{align*}
最后, 对于映射 $f:A\to B$ 和它的定义域 $A$ 的一个子集 $A'$, $f$ 在 $A'$ 上的限制, 记为 $f|_{A'}$, 是从 $A'$ 到 $B$ 的映射, 它将 $f(x)$ 指派给每个 $x\in A'$.

\end{document}